\documentclass[10pt,a4paper]{amsart}
\usepackage[margin=19mm]{geometry}
\usepackage{amsmath,amssymb,mathtools}
\usepackage[scr=rsfs,cal=euler]{mathalfa}
\usepackage{xfrac}
\usepackage[unicode,hidelinks,bookmarks=false]{hyperref}
\newcommand{\ie}{i.e.\ }
\newcommand{\cf}{cf.\ }
\newcommand{\resp}{respectively\ }
\newcommand{\Subsection}{\subsection}
\providecommand{\nobreakdash}{\nobreak\hbox{-}}
\setlength{\emergencystretch}{2em}
\multlinegap0pt
\usepackage{amssymb}
\newtheorem{proposition}{Proposition}
\newcommand{\lcm}{\operatorname{lcm}}
\title{Two Questions on $G$-Harmonic Tuples}
\author{Murali Menon}
\date{Source: arXiv:2608.15873v1 (CC BY 4.0)}
\begin{document}
\maketitle

\noindent\emph{Source and licence.} Two Questions on $G$-Harmonic Tuples. Version arXiv:2608.15873v1 (CC BY 4.0), distributed by arXiv under the Creative Commons Attribution 4.0 International licence, http://creativecommons.org/licenses/by/4.0/, as recorded in the arXiv licence metadata for this exact version and shown on the abstract page https://arxiv.org/abs/2608.15873v1. Full licence notice: \texttt{licenses/arxiv-2608.15873-CC-BY-4.0.txt}.\par\smallskip

\noindent\textit{Editorial scope.} Counted unit: the article's proposition prop:triple (source0.tex lines 173--178). The article's complete original proof of it is delivered with it (source0.tex lines 180--214, one \texttt{proof} environment of 35 lines). No mathematics is written, repaired or re-derived: the statement and the proof are the article's own lines, in the article's own notation, and no retained line is substituted or re-flowed.  No further source-local context is needed: every cross-reference of every retained line resolves inside the two delivered spans.  Nothing was omitted from any retained span, so the package carries no omission record.  No retained line contains a citation, so the wrapper carries no bibliography and no bibliography entry.  No retained line contains a figure environment, an image-inclusion command or drawing code, so no image is delivered and the package declares no asset dependency.  Editorial formatting only, in addition: an AMS \texttt{amsart} preamble that restates the article's own declarations and macros used by the retained lines, namely the environments \texttt{proposition} on separate counters, together with the standard packages those lines need.  The retained environments print in this document as Proposition 1 in order of appearance, whereas the article prints the same items with the numbering of its own full document.  The complete native source of the article as distributed in the arXiv e-print for this exact version is bundled unchanged as \texttt{sources/207790/source0.tex}.
\begin{proposition}\label{prop:triple}
Let $U_3,U_4,U_5$ be subgroups of a group $G$ with $[G:U_j]=6r_j$ for $3\le j\le 5$, where $r_3,r_4,r_5$ are pairwise coprime and $3\nmid r_3r_4r_5$. Then it is impossible that
\[
\alpha(U_3,U_4)=3,\qquad \alpha(U_3,U_5)=\alpha(U_4,U_5)=1.
\]
\end{proposition}

\begin{proof}
Suppose these values hold. Since the $r_j$ are pairwise coprime, $\gcd(6r_i,6r_j)=6$ and hence $\lcm(6r_i,6r_j)=6r_ir_j$ for distinct $i,j$. Therefore
\[
[G:U_3\cap U_4]=18r_3r_4,\qquad
[G:U_3\cap U_5]=6r_3r_5,\qquad
[G:U_4\cap U_5]=6r_4r_5,
\]
and the product formula gives
\[
|U_3U_4|=\frac{|G|}{2},\qquad |U_3U_5|=|U_4U_5|=\frac{|G|}{6}.
\]

Let $K=[G:U_3\cap U_4\cap U_5]$. Since $U_3\cap U_4\cap U_5$ lies in each pairwise intersection, $K$ is a common multiple of $18r_3r_4$, $6r_3r_5$ and $6r_4r_5$. We claim
\[
\lcm(18r_3r_4,\ 6r_3r_5,\ 6r_4r_5)=18r_3r_4r_5.
\]
Indeed, write $v_p$ for the $p$-adic valuation. At $p=3$ the term $18r_3r_4$ has $v_3=2$ (since $3\nmid r_3r_4r_5$) while $6r_3r_5$ and $6r_4r_5$ have $v_3=1$, so the maximum is $2=v_3(18r_3r_4r_5)$. At $p=2$, by pairwise coprimality at most one $r_j$ is even; the two terms containing that $r_j$ both have $2$-adic valuation $1+v_2(r_j)$ and the third has valuation $1$, so the maximum is $1+v_2(r_j)=v_2(18r_3r_4r_5)$ (and equals $1$ if every $r_j$ is odd). At any other prime $p$, at most one $r_j$ is divisible by $p$, and it occurs in two of the three terms, so the maximum is $v_p(r_j)=v_p(18r_3r_4r_5)$. Thus the claimed lcm holds, and in particular $18r_3r_4r_5 \mid K$, so we may write $K=18r_3r_4r_5\,m$ with $m\ge 1$.

By the product formula,
\[
|U_3(U_4\cap U_5)|
=\frac{|U_3|\,|U_4\cap U_5|}{|U_3\cap U_4\cap U_5|}
=\frac{\dfrac{|G|}{6r_3}\cdot\dfrac{|G|}{6r_4r_5}}{\dfrac{|G|}{18r_3r_4r_5\,m}}
=\frac{m}{2}\,|G|.
\]
Because $U_4\cap U_5\subseteq U_4$, we have $U_3(U_4\cap U_5)\subseteq U_3U_4$, so
\[
\frac{m}{2}|G| \le |U_3U_4|=\frac{|G|}{2},
\]
forcing $m=1$. Then $|U_3(U_4\cap U_5)|=|U_3U_4|$, and since $U_3(U_4\cap U_5)\subseteq U_3U_4$ are finite sets of equal cardinality (no claim that either is a subgroup is needed), they coincide: $U_3(U_4\cap U_5)=U_3U_4$. But $U_4\cap U_5\subseteq U_5$ gives $U_3(U_4\cap U_5)\subseteq U_3U_5$, whence
\[
\frac{|G|}{2}=|U_3U_4|=|U_3(U_4\cap U_5)|\le |U_3U_5|=\frac{|G|}{6},
\]
a contradiction.
\end{proof}

\end{document}
